\section{Preliminaries}

In this section we recall the basic semantic frameworks that have been used to interpret the intuitionistic modal language. We begin by fixing the language of $\IK$, which extends intuitionistic propositional logic with the modal operators $\Box$ and $\Diamond$. We then present the hierarchical semantics, in which modalities are interpreted over families of Kripke models organised into layers reflecting the intuitionistic growth of information. Finally, we introduce the bi-relational semantics, which equips a single Kripke model with both an intuitionistic epistemic relation and a modal accessibility relation, together with suitable interaction conditions. These definitions have as their main purpose to introduce some notation and terminology that will be used throughout.

\begin{definition}
    Let $\props$ be a countably infinite set of propositional symbols.
    The set of \defn{formulas} of \iml over $\props$ is defined by the following grammar:
    \[
        \varphi, \psi \Coloneqq
            \bot \mid
            p \mid
            \varphi \land \psi \mid
            \varphi \lor \psi \mid
            \varphi \to \psi    \mid
            \Diamond \varphi \mid
            \Box \varphi. 
    \]
    In this grammar, $p$ is any proposition symbol in $\props$.
    We use $\lnot \varphi$ as an abbreviation for $\varphi \to \bot$.
\end{definition}

Having introduced the language of $\IK$, we proceed to introduce its hierarchical semantics.

\begin{definition}\label[definition]{def:kmodel}
    A \defn{(relational) Kripke model} is a triple $\kmodel = \tuple{\Wkm, {\Rkm}, \Vkm}$ where:
        \begin{enumerate}
            \setlength{\itemsep}{0pt}
            \item $\Wkm$ is a non-empty set of \defn{worlds};
            \item $\Rkm$ is a binary \defn{accessibility relation} on $\Wkm$;
            \item $\Vkm\colon \Wkm \to 2^{\props}$ is a \defn{valuation function}.
        \end{enumerate}
        We omit the index $\kmodel$ from $\Wkm$, $\Rkm$, and $\Vkm$, when it is clear from the context.
        By a \defn{pointed Kripke model}, we mean a pair $\pair \kmodel s$ where $\kmodel$ is a Kripke model and $s \in \Wkm$. We often omit the parentheses and write simply $\npair\kmodel s$ when there is no ambiguity.
\end{definition}

% \lutz{I think the order is weird. 2.3 and 2.4 should come after 2.6}

\begin{figure}[t]
    \centering
    \footnotesize
    \begin{tikzpicture}[modal,node distance = 10mm]
        \node (c) {(b)};
        \node[point] (a) [above left=of c,label=left:$t$] {};
        \node[point] (b) [above right=of c,label=right:$v$] {};
        \node[point] (d) [below left=of c,label=left:$s$] {};
        \node[point] (e) [below right=of c,label=right:$u$] {};
        \path[->] (a) edge[dotted] node[above] {$\R$} (b);
        \path[->] (d) edge node[above] {$\R$} (e);
        \path[->] (d) edge node[left] {$\preccurlyeq$} (a);
        \path[->] (e) edge[dotted] node[right] {$\preccurlyeq$} (b);
    \end{tikzpicture}
    \hspace{1.5cm}
    \begin{tikzpicture}[modal,node distance = 10mm]
        \node (c) {(c)};
        \node[point] (a) [above left=of c,label=left:$v$] {};
        \node[point] (b) [above right=of c,label=right:$u$] {};
        \node[point] (d) [below left=of c,label=left:$s$] {};
        \node[point] (e) [below right=of c,label=right:$t$] {};
        \path[->] (a) edge[dotted] node[above] {$\R$} (b);
        \path[->] (d) edge node[above] {$\R$} (e);
        \path[->] (d) edge[dotted] node[left] {$\preccurlyeq$} (a);
        \path[->] (e) edge node[right] {$\preccurlyeq$} (b);
    \end{tikzpicture}
    \caption{Frame Conditions on Bi-relational Kripke models}
    \label{fig:bframe}
\end{figure}

\begin{definition}\label[definition]{def:imodel}
    A \defn{hierarchical model} is a pair $\hmodel = \tuple{\Khm, \Hleq}$ where:
    % \begin{enumerate}
    %     \setlength{\itemsep}{0pt}
    %     \item $\Khm$ is a non-empty set of Kripke models;
    %     \item $\Hleq$ is a preorder on $\Khm$.
    % \end{enumerate}
    $\Khm$ is a non-empty set of Kripke models; and $\Hleq$ is a preorder on $\Khm$ such that for all $\kmodel_1,\kmodel_2\in\Khm$ with $\kmodel_1 \Hleq \kmodel_2$:
        (a) ${\Wkm[1] \subseteq \Wkm[2]}$;
        (b) ${\Rkm[1] \subseteq \Rkm[2]}$; and
        (c) for all $s \in \Wkm[1]$, $\Vkm[1](s) \subseteq \Vkm[2](w)$.
    % \begin{enumerate}[label=(\alph*)]
    %     \setlength{\itemsep}{0pt}
    %     \item ${\Wkm[1] \subseteq \Wkm[2]}$;
    %     \item ${\Rkm[1] \subseteq \Rkm[2]}$;
    %     \item for all $s \in \Wkm[1]$, $\Vkm[1](s) \subseteq \Vkm[2](w)$.
    % \end{enumerate}
    %%We use $\kmodels_{\hmodel}$, and $\preccurlyeq_{\hmodel}$, to indicate $\kmodels$, and $\preccurlyeq$, respectively.
\end{definition}

% \begin{proposition}
%     If $\bmodel$ is a hierarchical model, then, $\preccurlyeq_{\bmodel}$ is a partial order.
% \end{proposition}
% \begin{proof}
%     We only need to prove that $\preccurlyeq_{\bmodel}$ is antisymmetric, i.e., if $\kmodel_1 \preccurlyeq_{\bmodel} \kmodel_2$ and $\kmodel_2 \preccurlyeq_{\bmodel} \kmodel_1$, then, $\kmodel_1 = \kmodel_2$.
%     The result is immediate from conditions (a)--(c).
% \end{proof}

\begin{definition}\label[definition]{def:sat:imodel}
    Let $\hmodel$ be a hierarchical model, $\kmodel \in \Khm$, and $s \in \Wkm$.
    Moreover, let $A$ be a formula.
    The \defn{forcing relation} $\hforces {\pair \kmodel s} A$ is defined inductively as:
    \[\begin{array}{lcl}
        \hforces {\pair \kmodel s} \bot
            & \text{iff} & \text{false} \\
        \hforces {\pair \kmodel s} p
            & \text{iff} & p \in \Vkm(s) \\
        \hforces {\pair \kmodel s} {B \land C}
            & \text{iff} & \hforces {\pair \kmodel s} B \text{ and } \hforces {\pair \kmodel s} C \\
        \hforces {\pair \kmodel s} {B \lor C}
            & \text{iff} & \hforces {\pair \kmodel s} B \text{ or } \hforces {\pair \kmodel s} C \\
        \hforces {\pair \kmodel s} {B \to C}
            & \text{iff} & \text{for all } \kmodel' \in \Khm,~
                \text{if } \kmodel \Hleq \kmodel' \text{ and } \hforces {\pair {\kmodel'} s} B
                    \text{, then, } \hforces {\pair {\kmodel'} s} C \\
        \hforces {\pair \kmodel s} \Diamond B & \text{iff}
            & \text{exists } t \in \Wkm \text{ such that } s \Rkm t \text{ and } \hforces {\pair \kmodel t} B \\
        \hforces {\pair \kmodel s} \Box B & \text{iff}
            & \text{for all } \kmodel' \in \Khm,~
                \text{if } \kmodel \Hleq \kmodel' \text{, then, }
                    \text{for all } t \in \Wof{\kmodel'},
                        \text{ if } s \Rof{\kmodel'} t \text{, then, } \hforces {\pair {\kmodel'} t} B.
    \end{array}\]
    We use $\nhforces {\pair \kmodel s} A$ to indicate the negation of $\hforces {\pair \kmodel s} A$.
    Immediately, $\hforces {\pair \kmodel s} \lnot B$ iff for all $\kmodel' \in \Khm$, if $\kmodel \Hleq \kmodel'$, then, $\nhforces {\pair {\kmodel'} s} B$.
    % We write $\forces\hmodel A$ iff $\hforces{\pair \kmodel s} A$ for all $\kmodel\in\Khm$ and $s\in \Wkm$.
    % Finally, we write $\Vdash A$, and call $A$ a tautology, iff $\Vdash^{\hmodel} A$ for all hierarchical models $\hmodel$.\lutz{we can't define tautology twice!}
\end{definition}

The hierarchical semantics aligns closely with the standard semantic treatment of intuitionistic propositional and first-order logic, where models are organised into layers ordered by inclusion, reflecting the monotonic growth of information \cite{Troelstra1988}. This conceptual proximity is one of its main attractions. This same feature, however, also leads to a significant limitation: the hierarchical semantics fails to enjoy the finite model property. A classical counterexample, originating in the work of Ono and Suzuki~\cite{ono1988}, is the formula $\Box \neg \neg \varphi \to \neg \neg \Box \varphi$. This formula can be falsified only in infinite hierarchical models.

Alternative to the hierarchical semantics for $\IK$ is the bi-relational semantics. We introduce this alternative semantics next.

\begin{definition}\label[definition]{def:bmodel}
    A \defn{birelational Kripke model} is a cuadruple of the form $\bmodel = \tuple{\Wbm, \Fbm, \Rbm, \Vbm}$ where:
    % \begin{enumerate}
    %     \setlength{\itemsep}{0pt}
    %     \item $\Wbm$ is a non-empty set of \defn{worlds};
    %     \item $\Fbm$, the \defn{future relation},    is a preorder (reflexive and transitive) on $\Wbm$;
    %     \item $\Rbm$ is a binary \defn{accessibility relation} on $\Wbm$;
    %     \item $\Vbm\colon \props \to 2^{\Wbm}$ is a \defn{valuation function}.
    % \end{enumerate}
    $\tuple{\Wbm, \Rbm, \Vbm}$ is a Kripke model, and 
    $\Fbm$, the \defn{future relation}, is a preorder on $\Wbm$
    such that
%%    Bi-relational Kripke models further satisfy the following conditions:
    \begin{enumerate}[label=(\alph*)]
        \setlength{\itemsep}{0pt}
        \item for all $s,t\in\Wbm$, if $s \Fbm t$ and $p \in \Vbm(s)$, then, $p \in \Vbm(t)$.
        \item for all $s,t,u\in\Wbm$, if $s \Fbm t$ and $s \Rbm u$, then, there is $v \in \Wbm$ such that $t \Rbm v$ and $u \Fbm v$;
        \item for all $s,t,u\in\Wbm$, if $s \Rbm t$ and $t \Fbm u$, then, there is $v \in \Wbm$ such that $s \Fbm v$ and $v \Rbm u$;
    \end{enumerate}
    The condition (a) above is referred to as \emph{monotonicity of valuations}.
    The conditions (b) and (c) above, depicted in \Cref{fig:bframe}, are commutativity conditions on future and accessibility relations.
    The conditions (b) and (c) are discussed in \cite[p.~51]{simpson1994}, under the names (F1) and (F2), respectively.
    % By a \defn{pointed birelational Kripke model}, we mean a pair $\pair \bmodel s$ where $\bmodel$ is a birelational Kripke model and $s \in \Wbm$.
    % We often omit the parentheses and write simply $\npair \bmodel s$ when there is no ambiguity.
    % If $v$ in (a) and (b) is unique, $\bmodel$ is called \defn{universal}.
    % We write $\preccurlyeq^{\equiv}$ and $\R^{\equiv}$ for the equivalence closure of $\preccurlyeq$ and $\R$, respectively.\lutz{what does that mean?}
    % We call $\bmodel$ \defn{cartesian} iff whenever $s \preccurlyeq^{\equiv} t$ and $t \R^{\equiv} u$, then, $s = u$.\todo{add index}
    % We use $\W_{\bmodel}$, $\preccurlyeq_{\bmodel}$, $\R_{\bmodel}$, and $\V_{\bmodel}$, to indicate $\W$, $\preccurlyeq$, $\R$, and $\V$, respectively.\lutz{We omit the index $\bmsymb$ ...}
\end{definition}

\begin{definition}\label[definition]{def:bsat}
    Let $\bmodel$ be a bi-relational model, $s \in \W_{\bmodel}$, and $A$ be a formula.
    The \defn{forcing relation} $\bforces s A$ is defined inductively as:
    \[\begin{array}{lcl}
            \bforces s \bot & \text{iff} & \text{false} \\
            \bforces s p & \text{iff} & p \in \V_{\bmodel}(s) \\
            %\bforces s \lnot B & \text{iff} & \text{for all } s \preccurlyeq_{\bmodel} t,~ \bmodel, t \nVdash B \\
            \bforces s {B \land C} & \text{iff} & \bforces s B \text{ and } \bforces s C \\
            \bforces s {B \lor C} & \text{iff} & \bforces s B \text{ or } \bforces s C \\
            \bforces s {B \to C} & \text{iff} & \text{for all } t \in \Wbm, \text{ if } s \Fbm t \text{ and } \bforces t B, \text{ then, } \bforces t C \\
            \bforces s {\Diamond B} & \text{iff} & \text{exists } s \in \Wbm \text{ such that } s \Rbm t \text{ and } \bforces t B \\
            \bforces s {\Box B} & \text{iff} & \text{for all } t,u \in \Wbm, \text{ if } s \Fbm t \text{ and } t \Rbm u, \text{ then, } \bforces u B.
    \end{array}\]
    We use $\nbforces s A$ to indicate the negation of $\bforces s A$.
    Immediately, $\bforces s \lnot B$ iff for all $t \in \Wbm$, if $s \Fbm t$, then, $\nbforces t B$.
    %We write $\forces \bmodel A$ iff $\bforces s A$ for all $s \in \Wbm$.
    % Moreover, we write $\sforce A$, and call $A$ a \defn{tautology}, iff $\bmodel \Vdash A$ for all bi-relational Kripke models $\bmodel$.
\end{definition}

We are now in position to comment on two key features of the bi-relational semantics for intuitionistic modal logic.
First, the clauses for $\Diamond$ and $\Box$ in \Cref{def:bsat} closely parallel those for the hierarchical semantics, and so for the existential and universal quantifiers in intuitionistic predicate logic~\cite{Troelstra1988}. This parallel reflects the constructive nature of the logic. Second,    conditions (b) and (c) in \Cref{def:bmodel} are traced back to the work of Plotkin in~\cite{PlotkinS86}, where they play a foundational role. Formally, condition (b) ensures the monotonicity of $\Diamond$: if $\bforces s {\Diamond A}$ and $s \Fbm t$, then $\bforces t {\Diamond A}$. In turn, condition (c) is motivated by the intuition that the formula $\lnot \Diamond A \to \Box \lnot A$ should hold---and by the fact that this implication fails if the condition is omitted.
While condition (b) can easily be regarded as natural, condition (c) is more debatable. This is also mentioned in~\cite{PlotkinS86,simpson1994}, where the authors themselves express greater confidence in the necessity of (b), acknowledging that (c) may be less well-justified.
In any case, condition (c) is nowadays accepted as inherent in a bi-relational model.

The following two propositions indicate the close semantic correspondence between the hierarchical and bi-relational models.

\begin{proposition}[{\cite[Lemma~8.1.2]{simpson1994}}]\label[proposition]{prop:hs2bks}
    Let $\hmodel$ be a hierarchical model.
    Define $\bmofhm = \tuple{\W, \Fsymb, \R, \V}$ where:
    \begin{enumerate}
        \setlength{\itemsep}{0pt}
        \item $\W = \setof{\pair\kmodel s}{ \kmodel \in \Khm \text{ and } s \in \Wkm}$;
        \item ${\Fsymb} = \setof{ (\pair{\kmodel_1}s, \pair{\kmodel_2} s) }{ \kmodel_1 \Hleq\kmodel_2 \text{ and } s\in    \Wkm[1]}$;
        \item $\R    = \setof{ (\pair{\kmodel}{s}, \pair{\kmodel}{t}) }{ \kmodel\in \Khm \text{ and } s \Rkm t}$;
        \item %$\V_{\hmodel}: \props \to 2^{\S_{\hmodel}}$ is s.t.\
            for all ${\pair \kmodel s} \in \W$, $\V(\npair \kmodel s) = \Vkm(s)$.
    \end{enumerate}
    It follows that:
        (a) $\bmofhm$ is a birelational Kripke model; and
        (b) for all formulas $A$, it follows that $\hforces{\pair\kmodel s} A$ ~iff~ $\xforces{\bmofhm}{\pair\kmodel s} A$.
    % \begin{enumerate}[label=(\alph*)]
    %     \setlength{\itemsep}{0pt}
    %     \item $\bmofhm$ is a birelational Kripke model;
    %     \item $\hforces{\pair\kmodel s} A$ ~iff~ $\xforces{\bmofhm}{\pair\kmodel s} A$; and
    %     % \item if $\hmodel$ is finite, then, $\bmodel_{\hmodel}$ is finite.
    % \end{enumerate}
\end{proposition}

\begin{proposition}\label[proposition]{prop:iml}
    Let $A$ be any formula.
    The following are equivalent:
    \begin{enumerate}[label=(\alph*)]
        \setlength{\itemsep}{0pt}
        \item for all hierarchical models $\hmodel$, $\kmodel \in \Khm$, and $s \in \Wkm$, it follows that $\hforces{\pair \kmodel s}{A}$;
        \item for all bi-relational models $\bmodel$, and $s \in \Wbm$, it follows that $\bforces{s}{A}$.
    \end{enumerate}
\end{proposition}

% In \cite[$\S$~8.1.2]{brauner11}, it is claimed that, in general, bi-relational Kripke models cannot be translated into hierarchical models.
% The stated reason is that there are more bi-relational Kripke models than intuitionistic modal models, though no formal justification for this claim is provided.
% At the same time, in \cite[p.~152]{simpson1994}, it is shown that cartesian bi-relational Kripke models can be translated into hierarchical models in a semantic preserving way.
% \valentin{When we started to look into this problem Carlos messaged Torben about the claim he makes in his book. \lutz{what claim?} He replied, but did not provide a good justification for what he says. From the result in Simpson's PhD thesis, I think it should be clear that the translation from cartesian bi-relational Kripke models to hierarchical models cannot preserve finiteness, or cartesian bi-relational Kripke models cannot be used to characterize $\IK$.}\lutz{I think finiteness is easily preserved. but cartesian models are too restriced to have the FMP}

The result in~\Cref{prop:hs2bks} establishes that every hierarchical model can be translated into a bi-relational Kripke model in a truth-preserving way. In turn, the result in \Cref{prop:iml} establishes that the two frameworks validate the same formulas. As pointed out in \cite{simpson1994}, the latter result is surprising given the structural different between the frameworks. At the same time, the model-theoretic behaviour of the two frameworks differs in a crucial respect. Bi-relational models enjoy the finite model property, whereas hierarchical models do not. This asymmetry provides a clear indication that, although hierarchical models can be faithfully encoded into bi-relational ones, the converse translation cannot in general be done in finiteness-preserving way. In this sense, the bi-relational semantics offers a strictly more robust model-theoretic setting, even though both frameworks validate the same intuitionistic modal formulas.
